\documentclass[11pt]{article}

\usepackage[margin=1in]{geometry}
\usepackage{amsmath}
\usepackage{amssymb}

\title{Notes on 0PN, 1PN, and TaylorF2 Chirp Parametrisations}
\author{}
\date{}

\begin{document}

\maketitle

\section{TaylorF2 phase convention}

The TaylorF2 waveform is a frequency-domain stationary-phase approximation.
For a nonspinning binary, its phase can be written schematically as
\begin{equation}
    \Psi(f)
    = 2\pi f t_c - \phi_c - \frac{\pi}{4}
    + \frac{3}{128\eta} v^{-5}
      \sum_{k=0}^{7}
      \left[a_k + b_k \log v\right] v^k ,
    \label{eq:taylorf2_phase}
\end{equation}
where
\begin{equation}
    v = (\pi M f)^{1/3}, \qquad
    M = m_1 + m_2, \qquad
    \eta = \frac{m_1 m_2}{(m_1 + m_2)^2}.
\end{equation}
The index $k$ corresponds to PN order $k/2$. Thus terms through $k=7$
correspond to phase corrections through 3.5PN. In the LAL TaylorF2 call used
here, with no explicit phase-order override and zero spins, the default
nonspinning TaylorF2 phasing includes nonzero terms up to 3.5PN.

\section{The 1PN TaylorF2 phase term}

Keeping only the Newtonian and 1PN TaylorF2 phase terms gives
\begin{equation}
    \Psi_{\mathrm{0PN+1PN}}(f)
    =
    2\pi f t_c - \phi_c - \frac{\pi}{4}
    + \frac{3}{128\eta} v^{-5}
      \left[
        1 + \alpha_2 v^2
      \right],
    \label{eq:taylorf2_1pn}
\end{equation}
with
\begin{equation}
    \alpha_2 = \frac{3715}{756} + \frac{55}{9}\eta .
\end{equation}
For an equal-mass binary, $\eta = 1/4$, so
\begin{equation}
    \alpha_2 = \frac{3715}{756} + \frac{55}{36}.
\end{equation}

\section{Factoring out frequency}

Up to the constant phase offset $-\phi_c - \pi/4$, the 0PN+1PN phase can be
written as
\begin{equation}
    \Psi_{\mathrm{0PN+1PN}}(f)
    =
    \mathrm{const}
    + 2\pi f t_c
    + A f^{-5/3}
    + B f^{-1},
\end{equation}
where
\begin{align}
    A &= \frac{3}{128\eta}(\pi M)^{-5/3}, \\
    B &= \frac{3}{128\eta}\alpha_2(\pi M)^{-1}.
\end{align}
Therefore one may factor out $2\pi f$ and define an effective
frequency-dependent time coordinate,
\begin{equation}
    \Psi_{\mathrm{0PN+1PN}}(f)
    =
    \mathrm{const}
    + 2\pi f \, \tau_{\mathrm{F2}}(f),
\end{equation}
where
\begin{equation}
    \tau_{\mathrm{F2}}(f)
    =
    t_c
    + \frac{A}{2\pi} f^{-8/3}
    + \frac{B}{2\pi} f^{-2}.
\end{equation}

The same algebraic factorisation is possible through 3.5PN. In general,
\begin{equation}
    \Psi(f)
    =
    \mathrm{const}
    + 2\pi f\,\tau_{\mathrm{F2}}(f),
\end{equation}
with
\begin{equation}
    \tau_{\mathrm{F2}}(f)
    =
    t_c
    + \frac{3}{256\pi\eta}
      \sum_{k=0}^{7}
      \left[a_k + b_k \log\left((\pi M f)^{1/3}\right)\right]
      (\pi M)^{(k-5)/3}
      f^{(k-8)/3}.
    \label{eq:tau_f2_factorised}
\end{equation}
The corresponding powers of $f$ inside $\tau_{\mathrm{F2}}$ are
\begin{center}
\begin{tabular}{c|c|c}
PN order & $k$ & power in $\tau_{\mathrm{F2}}(f)$ \\
\hline
0PN   & 0 & $f^{-8/3}$ \\
0.5PN & 1 & $f^{-7/3}$ \\
1PN   & 2 & $f^{-2}$ \\
1.5PN & 3 & $f^{-5/3}$ \\
2PN   & 4 & $f^{-4/3}$ \\
2.5PN & 5 & $f^{-1}$ \\
3PN   & 6 & $f^{-2/3}$ \\
3.5PN & 7 & $f^{-1/3}$
\end{tabular}
\end{center}
Some PN orders also include logarithmic corrections.

This factorisation is an algebraic statement about the phase. For the physical
chirp trajectory, the relevant time-frequency relation is instead
\begin{equation}
    t(f) = \frac{1}{2\pi}\frac{d\Psi}{df}.
    \label{eq:stationary_phase_time}
\end{equation}
This derivative relation is what should be used when extracting a chirp
trajectory from TaylorF2.

\section{0PN beta parametrisation}

At Newtonian order, the chirp can be written explicitly as
\begin{equation}
    f(\beta,t)
    =
    f_0
    \left(
      1 - \frac{8}{3}\beta t
    \right)^{-3/8}.
    \label{eq:f_0pn_beta}
\end{equation}
Here $f_0$ is the frequency at reference time $t=0$, and
\begin{equation}
    \beta
    =
    \frac{96}{5}
    \pi^{8/3}
    \left(\frac{G M_c}{c^3}\right)^{5/3}
    f_0^{8/3},
    \label{eq:beta_definition}
\end{equation}
where
\begin{equation}
    M_c
    =
    \frac{(m_1m_2)^{3/5}}{(m_1 + m_2)^{1/5}}
\end{equation}
is the chirp mass. The corresponding phase evolution is
\begin{equation}
    \phi(\beta,t)
    =
    -\frac{6\pi}{5}
    f_0
    \frac{
      \left(
        1-\frac{8}{3}\beta t
      \right)^{5/8}
    }{\beta}
    + \phi_0 .
    \label{eq:phi_0pn_beta}
\end{equation}
An additive constant may be chosen so that $\phi(0)=\phi_0$ if desired.

\section{Can the same construction be done at 1PN?}

At 1PN, the frequency evolution is still separable, but it is no longer
described by the single clean 0PN beta form. The 1PN evolution may be written
as
\begin{equation}
    \frac{df}{dt}
    =
    \frac{96}{5}
    \pi^{8/3}
    \left(\frac{G M_c}{c^3}\right)^{5/3}
    f^{11/3}
    \left[
      1
      -
      \alpha
      \left(
        \pi \frac{G M}{c^3} f
      \right)^{2/3}
    \right],
    \label{eq:dfdt_1pn}
\end{equation}
where
\begin{equation}
    \alpha = \frac{743}{336} + \frac{11}{4}\eta .
\end{equation}
For equal masses, $\eta=1/4$, giving $\alpha = 487/168$.

Using the same $\beta$ as in Eq.~\eqref{eq:beta_definition}, define
\begin{equation}
    x = \frac{f}{f_0},
    \qquad
    \epsilon
    =
    \alpha
    \left(
      \pi \frac{G M}{c^3} f_0
    \right)^{2/3}.
\end{equation}
Then Eq.~\eqref{eq:dfdt_1pn} becomes
\begin{equation}
    \frac{df}{dt}
    =
    \beta f_0 x^{11/3}
    \left[
      1 - \epsilon x^{2/3}
    \right].
\end{equation}
For a generic mass ratio, $\epsilon$ is an additional independent parameter.
However, for the equal-mass case used in the numerical examples, $\eta=1/4$ and
\begin{equation}
    M = 2^{6/5} M_c .
\end{equation}
Therefore $\epsilon$ can be eliminated in favour of $\beta$ and $f_0$.  From
Eq.~\eqref{eq:beta_definition},
\begin{equation}
    \left(
      \frac{G M_c}{c^3}
    \right)^{5/3}
    =
    \frac{5\beta}{96\pi^{8/3}f_0^{8/3}},
\end{equation}
so that
\begin{equation}
    \pi \frac{G M}{c^3} f_0
    =
    \left(
      \frac{5\beta}{24\pi f_0}
    \right)^{3/5}.
\end{equation}
For equal masses, $\alpha=487/168$, and hence
\begin{equation}
    \epsilon(\beta, f_0)
    =
    \frac{487}{168}
    \left(
      \frac{5\beta}{24\pi f_0}
    \right)^{2/5}.
    \label{eq:epsilon_equal_mass}
\end{equation}
The equal-mass 1PN chirp can therefore be written as
\begin{equation}
    \frac{df}{dt}
    =
    \beta f_0 x^{11/3}
    \left[
      1
      -
      \frac{487}{168}
      \left(
        \frac{5\beta}{24\pi f_0}
      \right)^{2/5}
      x^{2/3}
    \right],
    \qquad
    x = \frac{f}{f_0}.
    \label{eq:dfdt_1pn_equal_mass_beta}
\end{equation}
Thus, if $f_0$ is fixed, the equal-mass 1PN chirp is still a one-parameter
family labelled by $\beta$.  If $f_0$ is allowed to vary, then $\beta$ alone is
not enough; the 1PN correction depends on the dimensionless combination
$\beta/f_0$.

A slightly cleaner equal-mass parametrisation is obtained by defining
\begin{equation}
    \lambda
    =
    \pi \frac{G M}{c^3} f_0
    =
    \left(
      \frac{5\beta}{24\pi f_0}
    \right)^{3/5}.
\end{equation}
Then
\begin{equation}
    \frac{d x}{d(f_0 t)}
    =
    \frac{24\pi}{5}
    \lambda^{5/3}
    x^{11/3}
    \left[
      1
      -
      \frac{487}{168}
      \lambda^{2/3}x^{2/3}
    \right],
    \qquad
    x=\frac{f}{f_0}.
    \label{eq:dfdt_1pn_equal_mass_lambda}
\end{equation}
This form shows explicitly that $f(t)=f_0 x(f_0 t;\lambda)$ for equal masses.

The time-frequency relation can be integrated analytically:
\begin{equation}
    t(f)
    =
    \frac{1}{\beta}
    \int_1^{f/f_0}
    \frac{x^{-11/3}}{1-\epsilon x^{2/3}}\,dx .
\end{equation}
With the substitution
\begin{equation}
    u = x^{-2/3},
\end{equation}
this becomes
\begin{equation}
    t(f)
    =
    \frac{3}{2\beta}
    \left[
      \frac{u^4}{4}
      + \frac{\epsilon u^3}{3}
      + \frac{\epsilon^2 u^2}{2}
      + \epsilon^3 u
      + \epsilon^4 \log|u-\epsilon|
    \right]_{u=(f/f_0)^{-2/3}}^{u=1}.
    \label{eq:t_of_f_1pn}
\end{equation}
This gives an analytic implicit relation $t=t(f)$. In contrast to the 0PN
case, solving explicitly for $f(t)$ does not produce a simple closed form like
Eq.~\eqref{eq:f_0pn_beta}. In practice one would either numerically invert
Eq.~\eqref{eq:t_of_f_1pn} or integrate the 1PN ODE directly.

The phase can also be written analytically as a function of frequency:
\begin{equation}
    \phi(f)
    =
    \phi_0
    +
    2\pi
    \int_{f_0}^{f}
    f'
    \frac{dt}{df'}\,df'.
\end{equation}
Again, this naturally gives $\phi(f)$ rather than a simple explicit
$\phi(t)$.

\end{document}
